\scanpage{229}
(15) nous donne une \uline{réciproque} pour la condition envisagée pour que $(D_i,\omega_i)$ soit de Picard-Fuchs : il f. et s. que l'on ait (7). Cela résulte du th. de De Rham sur les fibres $X_s$ ....

\uline{Calcul de} $\mu(a,b)$ par voie ``algébrique'', sans calcul d'intégrales. \inserted{C'est local sur $S$.} La formule (5) nous ramène à calculer $\mu(a,b)$ \inserted{au voisinage de $s_0$,} pour \deleted{$a,b$} $a(s_0)$ et $b(s_0)$ assez voisins. Cela nous permet alors de supposer que sur $X/S$ il y a une connexion intégrable. On peut \unsure{$=$|même} supposer, par un choix convenable, que les sections $a,b$ sont horizontales. Alors on trouve (14)
\[
\mu(a,b)=\varphi(b(s))-\varphi(a(s))
\]
où $\varphi$ est défini par (10).

Il y aurait lieu \unsure{d'expliciter|d'expliquer} aussi une formule explicite, sans supposer que $a(s)$ et $b(s)$ soient horizontales. Faisons le calcul au voisinage de $s_0$, et soient $a_0,b_0$ les
\margin{228}
